% Conclusion

Redistricting is inherently temporal—districts are redrawn every decade—yet most research analyzes it in isolated snapshots. By applying a consistent algorithmic methodology across 2000, 2010, and 2020 census data, we demonstrate that algorithmic redistricting provides not merely a one-time improvement but a durable foundation for fair representation.

Our findings reveal a stark divergence: algorithmic compactness remains stable at 0.41--0.44 Polsby-Popper across two decades, while enacted district compactness improved by 11\% (0.298$\to$0.331) as state legislatures adapted to legal compactness scrutiny, though still trailing algorithmic plans by 19--24\%. Redistricting commissions—adopted by six states post-2010—reverse this negative trend but still underperform algorithmic baselines by 48\%.

Geographic stability analysis shows that algorithmic redistricting adapts gracefully to demographic change: 61\% of districts maintained substantial boundary overlap (IoU $>$ 0.7) from 2010 to 2020, even as 13 states gained or lost congressional seats. This resolves the apparent tension between compactness (optimizing district shape) and responsiveness (adapting to population shifts).

As the nation approaches the 2030 census, the 20-year record establishes algorithmic redistricting as a proven, durable approach. Our findings suggest that states may benefit from adopting redistricting commissions that use algorithmic baselines as starting points, deviating only for legally mandated considerations (VRA compliance, communities of interest). The alternative—continued reliance on partisan mapmakers with increasingly sophisticated tools—will likely yield further erosion of geometric compactness and public trust.

The question is no longer whether algorithmic redistricting works, but whether policymakers will adopt it before districting technology optimized for partisan advantage outpaces reform efforts entirely.
